% !TEX root = ../../../MathModel.tex
% 负责人：罗财能。
\subsubsection{具体分析}
\label{sec:q3:analysis}

问题三在运输任务中加入通信保障：运输无人机从起飞到返航，必须持续接入固定网关，或通过已完成建链的中继无人机接入网关。
山区地形既影响航段巡航高度、爬升时间和能耗，也会遮挡通信视线，因此航线可飞并不意味着沿途通信可用。
即使起点和服务区均能通信，爬升、巡航或下降过程仍可能出现覆盖缺口，不能只检查端点。

本问的决策包括不可拆货箱的组批与机型选择、服务区访问顺序、运输开始时刻、实体无人机和共享电池分配，以及中继悬停位置、高度、服务时段和能源组件分配。
这些决策通过时间相互制约：运输任务必须等待中继就位建链，中继提前到达会增加悬停能耗，过晚到达则可能导致硬时限任务不可行；同一中继服务时段又可以同时保障多个满足链路条件的运输架次。
因此，将运输先独立排定、再事后补充中继，容易破坏原有时限或资源约束，需要进行联合调度。

本文采用“地形与链路预计算、多起点联合构造、固定结构时序优化、独立可行性核验”的求解流程。
首先在有限候选悬停位置上筛选具备回传和接入能力的中继点；随后同时安排运输与中继任务，并用约束规划重新分配开始时刻和实体资源；最终从原始输入重算货箱、时限、能耗、充电及整个飞行过程的通信覆盖。
题目要求综合权衡及时性、完成时间、能源和架次数，但未给出唯一权重或优先级。
为突出应急物资的及时交付，本文将全部硬约束作为可行性前提，再按加权延迟、联合完成时间、总能耗、总架次数的字典序选择方案。
该优先级是本文的建模选择，不将有限候选搜索的结果表述为原问题的全局最优解。
